%-------------------------
% Resume in Latex
% Author : Jake Gutierrez
% Based off of: https://github.com/sb2nov/resume
% License : MIT
%------------------------

\documentclass[letterpaper,11pt]{article}

\usepackage{latexsym}
\usepackage[empty]{fullpage}
\usepackage{titlesec}
\usepackage{marvosym}
\usepackage[usenames,dvipsnames]{color}
\usepackage{verbatim}
\usepackage{enumitem}
\usepackage[hidelinks]{hyperref}
\usepackage{fancyhdr}
\usepackage[english]{babel}
\usepackage{tabularx}
\usepackage{ragged2e}
\usepackage{microtype}
% pdfLaTeX Unicode mappings; XeLaTeX handles Unicode directly.
\ifdefined\pdfglyphtounicode
  \input{glyphtounicode}
\fi


%----------FONT OPTIONS----------
% sans-serif
% \usepackage[sfdefault]{FiraSans}
% \usepackage[sfdefault]{roboto}
% \usepackage[sfdefault]{noto-sans}
% \usepackage[default]{sourcesanspro}

% serif
% \usepackage{CormorantGaramond}
% \usepackage{charter}


\pagestyle{fancy}
\fancyhf{} % clear all header and footer fields
\fancyfoot{}
\renewcommand{\headrulewidth}{0pt}
\renewcommand{\footrulewidth}{0pt}

% Adjust margins
\addtolength{\oddsidemargin}{-0.5in}
\addtolength{\evensidemargin}{-0.5in}
\addtolength{\textwidth}{1in}
\addtolength{\topmargin}{-.5in}
\addtolength{\textheight}{1.0in}

\urlstyle{same}

\raggedbottom
\raggedright
\setlength{\tabcolsep}{0in}

% Sections formatting
\titleformat{\section}{
  \vspace{-4pt}\scshape\raggedright\large
}{}{0em}{}[\color{black}\titlerule \vspace{-5pt}]

% Ensure that generate pdf is machine readable/ATS parsable
\ifdefined\pdfgentounicode
  \pdfgentounicode=1
\fi

%-------------------------
% Custom commands
\newcommand{\resumeItem}[1]{
  \item\small{
    {#1 \vspace{-2pt}}
  }
}

\newcommand{\resumeSubheading}[4]{
  \vspace{-2pt}\item
    \begin{tabular*}{0.97\textwidth}[t]{l@{\extracolsep{\fill}}r}
      \textbf{#1} & #2 \\
      \textit{\small#3} & \textit{\small #4} \\
    \end{tabular*}\vspace{-7pt}
}

\newcommand{\resumeSubSubheading}[2]{
    \item
    \begin{tabular*}{0.97\textwidth}{l@{\extracolsep{\fill}}r}
      \textit{\small#1} & \textit{\small #2} \\
    \end{tabular*}\vspace{-7pt}
}

\newcommand{\resumeProjectHeading}[2]{
    \item
    \begin{tabularx}{0.97\textwidth}{@{}>{\raggedright\arraybackslash}X r@{}}
      \small#1 & #2 \\
    \end{tabularx}\vspace{-7pt}
}

\newcommand{\resumeSubItem}[1]{\resumeItem{#1}\vspace{-4pt}}

\renewcommand\labelitemii{$\vcenter{\hbox{\tiny$\bullet$}}$}

\newcommand{\resumeSubHeadingListStart}{\begin{itemize}[leftmargin=0.15in, label={}]}
\newcommand{\resumeSubHeadingListEnd}{\end{itemize}}
\newcommand{\resumeItemListStart}{\begin{itemize}[rightmargin=0.15in]\justifying}
\newcommand{\resumeItemListEnd}{\end{itemize}\vspace{-5pt}}

%-------------------------------------------
%%%%%%  RESUME STARTS HERE  %%%%%%%%%%%%%%%%%%%%%%%%%%%%


\begin{document}

\begin{center}
    \textbf{\Huge \scshape Jin Tao} \\ \vspace{1pt}
    \small \href{mailto:j.tao@campus.tu-berlin.de}{\underline{j.tao@campus.tu-berlin.de}} $|$
    \href{https://www.linkedin.com/in/jin-tao-89035a289/}{\underline{linkedin.com/in/jin-tao-89035a289}} $|$
    \href{https://github.com/JinTTTT}{\underline{github.com/JinTTTT}}
    \\[6pt]
    {\large Selected Work Samples}
\end{center}

\section{Home Assistant Robot}
\noindent\small\textbf{Demo:} \href{https://youtu.be/3TWB9y0iEz0}{\underline{youtu.be/3TWB9y0iEz0}}
\resumeItemListStart
    \resumeItem{Lead development of a TurtleBot4 home assistant platform at ZEKI, integrating SLAM, navigation, autonomous exploration, semantic mapping, and an AI agent that translates user requests into robot actions.}
    \resumeItem{Implement a distributed onboard architecture: Raspberry Pi runs SLAM and navigation; NVIDIA Jetson Nano runs semantic mapping and the agent system.}
\resumeItemListEnd

\section{TaoRobot: Autonomous Driving Full Stack}
\noindent\small\textbf{Code and demos:} \href{https://github.com/JinTTTT/taorobot}{\underline{github.com/JinTTTT/taorobot}}
\resumeItemListStart
    \resumeItem{Sole developer of an open-source ROS 2 autonomy stack, owning architecture, algorithm implementation, and integration in Gazebo simulation.}
    \resumeItem{Implemented pose-graph SLAM, particle-filter localization, A* path planning with path smoothing, and pure pursuit trajectory tracking. The repository includes demonstrations of mapping, localization, and goal-directed navigation.}
\resumeItemListEnd

\section{Magnetic Joint-State Estimation}
\noindent\small\textbf{Code:} \href{https://github.com/JinTTTT/magnetic_field_simulation_for_rollPitchYaw_estimation}{\underline{GitHub repository: magnetic orientation estimation}}
\resumeItemListStart
    \resumeItem{Proposed and developed a ball-joint orientation estimation system using a magnet and two low-cost three-axis magnetic sensors to estimate yaw, pitch, and roll.}
    \resumeItem{The implementation combines magnetic-field simulation, physical-model calibration, and a lookup-based estimator. Evaluated the system against an Xsens MTi-630 IMU; the magnetic sensors provide the measurements used during normal operation.}
\resumeItemListEnd

\section{3D Object Pose Estimation with a Single RGB Camera}
\noindent\small\textbf{Code:} \href{https://github.com/JinTTTT/3d_obj_pose_estimation_with_only_rgb_camera}{\underline{GitHub repository: RGB-only object pose estimation}}
\resumeItemListStart
    \resumeItem{Developed an OpenCV-based project for estimating 3D object pose using a single RGB camera.}
\resumeItemListEnd

\end{document}
